\chapter{Core Theory: AI Compiler Optimization Principles}
\label{chap:ch13}
\typeout{START_CHAPTER "ch13" \theabspage}

What makes an AI compiler ``correct'' for production decode---pass ordering that looks elegant on a whiteboard, or schedules that measurably cut bytes/token and launches/token on the SKU you ship? Part~VI opens with \emph{Core Theory}: the optimization principles that sit \emph{below} Chapter~12 execution modes and \emph{above} dialect-specific backends (MLIR, Glow, TVM) \citep{chen2020compilerbasedhardw,yirage2026,taxbreak2026}. Every section binds IR transforms to residency on GPU, CPU, and NPU---not abstract big-O \citep{patel2025llminference,amd2024xdna,amd2024aie,dlbook}.

TaxBreak dense Llama-3.2-1B averages 847.5 kernel launches per output token on H100; OLMoE-1B/7B averages 9,305 \citep{taxbreak2026}. Memory-Floor reports $1.259\times$ median speedup on H100 from CUDA Graph replay alone at BS${=}1$, ctx${=}2048$---orchestration relief that tiling and fusion must extend into bytes, not replace \citep{memoryfloor2026,nvidia2024cudagraphs}. Compiler theory is how those counters move together: tile so data stays on-chip, fuse so boundaries disappear, layout so hardware reads coalesce, parallelize so batch-1 decode still saturates memory paths, and split backends so one IR lowers per target without forking model code \citep{yirage2026,chen2020compilerbasedhardw,semianalysis2024tma,dao2022flashattention,kwon2023vllm,liu2024deepseekv3,dao2023flashdecoding}.

This chapter is deliberately \emph{theory-first}: no MLIR syntax yet (Chapter~14), no manual CUDA (Part~VII). The goal is a shared vocabulary for pass designers---when someone proposes a new fusion rule or tile heuristic, Table~\ref{tab:ch13_pillars} names which counter must move and which prior chapter constraint applies \citep{yirage2026,chen2020compilerbasedhardw,patel2025llminference}. Industrial teams that skip this layer debate dialect patches for weeks while bytes/token flatlines because layout and tiling were never co-designed \citep{memoryfloor2026,taxbreak2026,kwon2023vllm}.

Fregly-style goodput analysis applies throughout: every transform is judged on batch-1 decode at representative context buckets, not on square GEMM microbenchmarks or prefill-only traces \citep{patel2025llminference,memoryfloor2026,taxbreak2026}. The five sections mirror how production compilers order concerns---tile shapes constrain fusion legality; layout determines whether fused regions can load coalesced data; parallel axes decide whether extra blocks help bandwidth; backend split decides which target owns each residency class \citep{yirage2026,chen2020compilerbasedhardw,amd2024xdna,amd2024aie,dlbook}.

Readers coming from classic optimizing-compiler courses should expect different objectives: decode LLM compilers optimize \emph{moved bytes} and \emph{launches} under numerics constraints, not loop nest locality alone \citep{taxbreak2026,memoryfloor2026,patel2025llminference,dao2022flashattention,yirage2026,chen2020compilerbasedhardw,dlbook}. When in doubt, re-run the Chapter~1 counter worksheet on the same model checkpoint before and after a pass change---theory without measurement is slideware \citep{patel2025llminference,taxbreak2026,memoryfloor2026,semianalysis2024tma,yirage2026,chen2020compilerbasedhardw,kwon2023vllm,liu2024deepseekv3,dao2022flashattention,dao2023flashdecoding,dlbook}.

\section{tiling theory}
\label{sec:ch13_tiling_theory}

% AUTO_VISUAL:fig_tiling_theory_pipeline

\begin{figure}[t]
\centering
\begin{tikzpicture}[vis base, scale=0.88, transform shape]
 \node[vis pipeline node] (s0) at (-5.03,0) {\footnotesize Global tensor\\HBM/DDR};
 \node[vis arch node] (s1) at (-1.68,0) {\footnotesize Tile scheduler\\block sizes};
 \node[vis pipeline node] (s2) at (1.67,0) {\footnotesize On-chip tile\\SMEM / L2 / AIE};
 \node[vis arch node] (s3) at (5.03,0) {\footnotesize Compute\\MMA / vector};
 \draw[vis arrow] (s0.east) -- (s1.west);
 \draw[vis arrow] (s1.east) -- (s2.west);
 \draw[vis arrow] (s2.east) -- (s3.west);
\end{tikzpicture}
\caption{Tiling theory: map global decode tensors into on-chip tiles before compute.}
\label{fig:tiling_theory_pipeline}
\end{figure}

\textbf{Tiling theory} asks which sub-blocks of weights, activations, and KV cache fit in registers, shared memory, L2, or NPU tile SRAM for each decode step \citep{dao2022flashattention,semianalysis2024tma,nvidia2022hopper}. FlashAttention is the canonical IO-aware tiling result for attention: recomputation and block-wise softmax avoid materializing $P$ in HBM \citep{dao2022flashattention,dao2023flashdecoding}. The compiler question is general---every matmul, norm, and MLP in decode needs a tile shape that respects bandwidth and capacity, not only MMA convenience \citep{patel2025llminference,memoryfloor2026,taxbreak2026}.

On Hopper, tiling couples to TMA: async copies move KV tiles while MMA consumes prior tiles \citep{nvidia2022hopper,semianalysis2024tma}. Wrong tile sizes leave TMA underutilized or spill partial outputs to HBM, showing up as bytes/token regression versus a smaller but legal tile \citep{memoryfloor2026,yirage2026}. On XDNA/AIE, tiling is static: DMA chains and tile slots are fixed at compile time; ``dynamic'' GPU tiles do not lower verbatim \citep{amd2024xdna,amd2024aie,chen2020compilerbasedhardw}. CPU decode tiling is cache blocking---NUMA-local KV and blocked GEMMs keep lines hot across norm--matmul chains \citep{dlbook,kwon2023vllm}.

\textbf{Analytic tile sizing.} A first-principles decode tile budget starts from hidden size $d$, head count $h$, KV dtype bytes $b_{kv}$, and context length $L$: bytes moved per token scale with reading prior KV plus writing new KV slots \citep{kwon2023vllm,dao2022flashattention,patel2025llminference}. Tile width along $L$ trades SMEM for trip count; tile width along $d$ trades MMA efficiency for reload width \citep{dao2023flashdecoding,semianalysis2024tma,memoryfloor2026}. Compilers should expose these trade-offs in IR annotations so autotuners search a \emph{legal} polytope rather than all factorizations of $L$ \citep{yirage2026,chen2020compilerbasedhardw,taxbreak2026,dlbook}. When analytic estimates disagree with profiles by more than a few percent, trust profiles---roofline models omit allocator, paging, and framework overhead \citep{memoryfloor2026,patel2025llminference,taxbreak2026}.

Compiler passes should emit tile metadata in IR (tile sizes, double-buffer depth, residency class) so backends reject illegal schedules before codegen \citep{yirage2026,chen2020compilerbasedhardw,liu2024deepseekv3}. Autotuners that search FLOPs-only tile shapes without bytes/token feedback recreate prefill wins and decode losses \citep{patel2025llminference,taxbreak2026,memoryfloor2026}. Profile each candidate tile on BS${=}1$ decode cells, not only square GEMM microbenchmarks \citep{taxbreak2026,semianalysis2024tma}.

\textbf{Prefill vs decode tiling.} Prefill benefits from large token-parallel tiles that amortize weight loads across many queries; decode at BS${=}1$ reuses weights once per step and instead fights KV reload along the sequence axis \citep{dao2022flashattention,dao2023flashdecoding,patel2025llminference}. A tile policy tuned on prefill shapes often selects block sizes that maximize MMA occupancy while increasing KV trips---exactly the regression Memory-Floor documents when bytes/token stay high despite faster attention kernels \citep{memoryfloor2026,taxbreak2026}. Compiler tile search must therefore tag phases explicitly: prefill buckets, decode buckets, and speculative-draft buckets when applicable \citep{kwon2023vllm,yirage2026}.

\textbf{Double buffering and pipeline depth.} Hopper pipelines pair TMA loads with MMA compute; tile theory includes how many stages fit in shared memory without spilling partial softmax or norm statistics \citep{nvidia2022hopper,semianalysis2024tma,dao2022flashattention}. Shallow pipelines underutilize memory bandwidth; deep pipelines overflow SMEM and force silent HBM spill \citep{memoryfloor2026,yirage2026}. On AIE, pipeline depth is fixed at compile time---there is no runtime knob to ``add another buffer'' when decode context grows \citep{amd2024xdna,amd2024aie,chen2020compilerbasedhardw}.

\textbf{Autotuning discipline.} TVM-style autotuners historically optimized for peak FLOPs on convolution and dense GEMM shapes \citep{dlbook,chen2020compilerbasedhardw}. LLM decode autotuning must swap the objective: minimize bytes/token subject to numerics and legality constraints from Chapters~10--11 \citep{dao2022flashattention,yirage2026,patel2025llminference}. Every autotune candidate should log tile dimensions \emph{and} measured HBM traffic per token---otherwise teams ship tiles that win on roofline FLOPs charts and lose on fleet p99 \citep{taxbreak2026,memoryfloor2026,semianalysis2024tma}.

\paragraph{Multi-hardware delta.} Identical FLOP count with different tile policies can differ by an order of magnitude in DDR traffic on XDNA versus TMA-backed Hopper tiles \citep{amd2024xdna,nvidia2022hopper,semianalysis2024tma,dlbook}.

\begin{example}
Decode cell (BS${=}1$, ctx${=}2048$): shrinking attention block size to fit shared memory may \emph{increase} loop iterations and bytes/token if KV reloads dominate---tile theory must co-optimize residency and trip count \citep{dao2022flashattention,memoryfloor2026,taxbreak2026,yirage2026}.
\end{example}

Industrial tile review checklists include: SMEM/TMA budget worksheet, KV bytes per trip, prefill vs decode bucket tags, and autotune objective function \citep{semianalysis2024tma,memoryfloor2026,yirage2026,chen2020compilerbasedhardw,patel2025llminference}. Teams that approve tile changes without decode profiles often ship kernels that win on nsys ``compute'' rows and lose on ``dram'' rows \citep{taxbreak2026,memoryfloor2026,dao2022flashattention,amd2024xdna,dlbook}.

\paragraph{Residency classes before tile shapes.}
The first tiling decision is not a number but a class: every tensor in the decode graph must be classified as register-resident, shared-memory or tile-SRAM resident, L2-resident, or global-memory streamed \citep{semianalysis2024tma,nvidia2022hopper,memoryfloor2026}. Only after classification does the pass choose concrete tile widths that satisfy the class boundaries; choosing MMA-preferred widths first and hoping residency follows is how fused graphs silently spill \citep{dao2022flashattention,yirage2026}. The classification itself is hardware-derived---Hopper exposes shared memory and L2 as separate budgets, XDNA fixes tile addresses at compile time, and CPU treats cache lines as the residency unit \citep{amd2024xdna,amd2024aie,dlbook}. IR should carry the class as an attribute so a later pass cannot demote a register tensor to shared memory without a legality review \citep{chen2020compilerbasedhardw,patel2025llminference}.

\paragraph{Tile sweep methodology.}
Tile autotuning becomes reproducible when the search is bounded by the residency classification first and measured on decode counters second \citep{patel2025llminference,taxbreak2026,memoryfloor2026}. Generate candidates by varying one axis at a time---KV-block width, head-block width, double-buffer depth---while holding the other axes at the analytic default; log bytes/token and ms/token for every candidate at the same context bucket \citep{dao2023flashdecoding,semianalysis2024tma}. Narrow the sweep around the best bytes/token point, not the best SM-utilization point, because batch-1 decode rewards fewer and wider memory transactions \citep{memoryfloor2026,taxbreak2026}. Keep the full sweep table in CI artifacts so a later chip-model update can re-run the same procedure without re-deriving the methodology \citep{yirage2026,chen2020compilerbasedhardw}.

\paragraph{Decode-specific tile invariants.}
Three invariants separate decode tiling from prefill tiling. First, weight tiles are read once per token, so they should stream through the largest reusable footprint rather than be re-tiled around a batch of queries \citep{kwon2023vllm,patel2025llminference}. Second, KV tiles are read on every new token, so their width along the sequence axis directly sets bytes/token and the shared-memory high-water \citep{dao2022flashattention,dao2023flashdecoding}. Third, softmax and norm statistics must remain resident across tile boundaries; a tile that invalidates their lifetime is illegal regardless of how fast it loads \citep{dao2022flashattention,yirage2026}. These invariants belong in the compiler test suite as negative tests: a pass that proposes an illegal decode tile should fail CI with the invariant name \citep{memoryfloor2026,chen2020compilerbasedhardw}.

\paragraph{Capacity ledger as review artifact.}
Every tiling PR should carry a capacity ledger: tensor names, assigned residency classes, byte estimates at the shipped context bucket, and the measured bytes/token from the decode profile \citep{semianalysis2024tma,memoryfloor2026,yirage2026}. The ledger turns an abstract tile change into a testable claim---if the ledger says an activation is shared-memory resident but the profile shows DRAM traffic growing with that tensor's size, the claim is falsified in review \citep{taxbreak2026,patel2025llminference}. Compiler-generated plans should emit the ledger automatically from IR metadata so it cannot go stale \citep{chen2020compilerbasedhardw,yirage2026}. Reviewers should reject tile PRs whose ledger and profile disagree, because that disagreement is the earliest symptom of silent spill \citep{memoryfloor2026,taxbreak2026}.

\section{fusion theory}
\label{sec:ch13_fusion_theory}

% AUTO_VISUAL:fig_fusion_theory_architecture

\begin{figure}[t]
\centering
\begin{tikzpicture}[vis base, scale=0.88, transform shape]
 \node[vis arch node] (s0) at (-5.03,0) {\footnotesize Op A\\};
 \node[vis arch node] (s1) at (-1.68,0) {\footnotesize Op B\\};
 \node[vis arch node] (s2) at (1.67,0) {\footnotesize Fused AB\\on-chip};
 \node[vis arch node] (s3) at (5.03,0) {\footnotesize HBM\\if illegal};
 \draw[vis arrow] (s0.east) -- (s1.west);
 \draw[vis arrow] (s1.east) -- (s2.west);
 \draw[vis arrow] (s2.east) -- (s3.west);
\end{tikzpicture}
\caption{Fusion theory: merge producer--consumer ops when residency allows; else spill.}
\label{fig:fusion_theory_architecture}
\end{figure}

\textbf{Fusion theory} governs when operator boundaries may disappear into MegaKernels or fused subgraphs \citep{yirage2026,dao2022flashattention,taxbreak2026}. Each unfused boundary exports activations to HBM; Chapter~1 counted the result as launches/token and bytes/token \citep{taxbreak2026,memoryfloor2026,patel2025llminference,kwon2023vllm}. Fusion is not ``always on''---it is constrained by register/shared-memory budgets, softmax carrier lifetimes (Chapter~10), and graph capture legality (Chapter~12) \citep{dao2022flashattention,nvidia2024cudagraphs,yirage2026,semianalysis2024tma}.

Greedy fusion (merge any adjacent ops) often illegalizes attention--MLP seams: online softmax state $(m,\ell,o)$ must survive inside the fused region or bytes explode silently \citep{dao2022flashattention,dao2023flashdecoding,memoryfloor2026}. Compiler cost models should penalize HBM round trips and launches, not only instruction count \citep{chen2020compilerbasedhardw,yirage2026,patel2025llminference,taxbreak2026}. YiRage-style legality checks encode Chapter~2 iron rules as IR predicates so fusion plans fail in compile time rather than in fleet A/B \citep{yirage2026,chen2020compilerbasedhardw,liu2024deepseekv3,amd2024aie}.

\textbf{Fusion boundaries as contracts.} Treat each candidate fusion edge as a contract between producer and consumer: the producer promises not to materialize to HBM; the consumer promises register/SMEM budget to consume in-place \citep{yirage2026,memoryfloor2026,dao2022flashattention}. When either side breaks the contract---for example LayerNorm statistics exported because the following GEMM needs a different layout---fusion must split \citep{taxbreak2026,patel2025llminference}. Graph capture (Chapter~12) records whatever fusion boundaries exist; it cannot invent new ones \citep{nvidia2024cudagraphs,memoryfloor2026}.

\textbf{Horizontal vs vertical fusion.} \emph{Horizontal} fusion merges operators at the same tensor rank (norm + linear + activation). \emph{Vertical} fusion merges stages along the decode pipeline (attention block + MLP block into one MegaKernel region) \citep{yirage2026,dao2022flashattention,taxbreak2026}. Vertical fusion delivers the largest bytes/token wins but hits legality first; horizontal fusion is easier but may leave launches/token high \citep{taxbreak2026,nvidia2024cudagraphs,patel2025llminference}. Compiler heuristics should attempt vertical plans only after horizontal plans pass carrier and layout checks \citep{yirage2026,chen2020compilerbasedhardw,semianalysis2024tma}.

\textbf{Fusion and execution modes.} Chapter~12 maps eager, graph, and MegaKernel modes to launch and byte counters; fusion theory explains \emph{which} boundaries each mode can erase \citep{nvidia2024cudagraphs,yirage2026,memoryfloor2026,taxbreak2026,patel2025llminference}. Graph capture records existing fusion; MegaKernel mode requires fusion plans that satisfy carrier legality \citep{dao2022flashattention,dao2023flashdecoding,yirage2026,chen2020compilerbasedhardw,semianalysis2024tma,nvidia2022hopper,amd2024aie,amd2024xdna,dlbook,kwon2023vllm,liu2024deepseekv3}. Teams stuck in graph mode with high bytes/token are blocked by fusion theory, not by CUDA Graph APIs \citep{memoryfloor2026,taxbreak2026,nvidia2024cudagraphs,patel2025llminference,yirage2026,chen2020compilerbasedhardw,dao2022flashattention,semianalysis2024tma,dlbook}. Document fusion-closed regions in runbooks so on-call engineers know whether to tune capture buckets or rewrite fusion plans \citep{yirage2026,chen2020compilerbasedhardw,patel2025llminference,taxbreak2026,memoryfloor2026,nvidia2024cudagraphs,amd2024xdna,amd2024aie,dlbook,kwon2023vllm,liu2024deepseekv3,dao2022flashattention,dao2023flashdecoding,nvidia2022hopper,semianalysis2024tma}.

\textbf{Cost model sketch.} A minimal fusion cost model scores each plan as $\alpha \cdot \mathrm{bytes} + \beta \cdot \mathrm{launches} + \gamma \cdot \mathrm{spill\_penalty}$, with spill penalty derived from residency metadata rather than FLOPs \citep{memoryfloor2026,taxbreak2026,patel2025llminference,yirage2026}. Tuning $(\alpha,\beta,\gamma)$ per SKU is acceptable; tuning without bytes/token telemetry is not \citep{semianalysis2024tma,dlbook}.

\paragraph{Multi-hardware delta.} Hopper may fuse QKV+softmax+proj when shared memory permits; XDNA may require separate static subgraphs linked by DMA; CPU fusion is often cache-scope fusion across norm--linear chains \citep{nvidia2022hopper,amd2024xdna,dlbook,chen2020compilerbasedhardw,kwon2023vllm}.

\begin{example}
Fusing LayerNorm with the following GEMM removes one full-width activation write/read per layer per token---often larger decode win than tuning MMA tile width alone \citep{memoryfloor2026,taxbreak2026,yirage2026,patel2025llminference}.
\end{example}

MoE layers add gather/scatter fusion constraints: expert routing may remain host-visible while MLP experts fuse internally---fusion theory must specify \emph{which} subgraph is fusion-closed \citep{taxbreak2026,liu2024deepseekv3,kwon2023vllm,yirage2026,nvidia2024cudagraphs,chen2020compilerbasedhardw,amd2024xdna,semianalysis2024tma,dlbook,dao2022flashattention,dao2023flashdecoding,nvidia2022hopper,amd2024aie,patel2025llminference,memoryfloor2026}.

Fusion regression tests should diff IR fusion boundaries against a golden eager graph: any new HBM edge between previously fused ops is a release blocker unless documented as intentional \citep{yirage2026,chen2020compilerbasedhardw,memoryfloor2026,taxbreak2026,patel2025llminference,dao2022flashattention,nvidia2024cudagraphs,semianalysis2024tma,dlbook}. This is how compiler teams prevent ``silent unfusion'' when dialect refactors reorder passes \citep{chen2020compilerbasedhardw,yirage2026,liu2024deepseekv3,amd2024aie,amd2024xdna}.

\paragraph{Fusion plans as IR contracts.}
A fusion plan should be representable as a first-class IR object: the fused region, the boundaries it removes, the buffers it keeps on-chip, and the numerical assumptions it makes \citep{yirage2026,chen2020compilerbasedhardw,dao2022flashattention}. Without that object, a pass that merges two ops and a pass that merely reorders them look identical in review, and silent unfusion can hide for releases \citep{memoryfloor2026,taxbreak2026}. The contract should be diffable in code review: a PR that changes fusion boundaries changes the contract diff, and the byte sheet must move with it \citep{patel2025llminference,chen2020compilerbasedhardw}. Compilers that encode contracts this way also enable rollback at the region level instead of at the whole-kernel level \citep{yirage2026,memoryfloor2026}.

\paragraph{Layout-changing producers.}
Many fusion opportunities die at layout conversions rather than at compute: a producer that emits a transposed or interleaved activation forces the consumer to either reshape in place or load strided \citep{semianalysis2024tma,dlbook,patel2025llminference}. Fusion theory must therefore include a layout negotiation step where producer and consumer agree on one layout before the boundary disappears \citep{chen2020compilerbasedhardw,yirage2026}. When no common layout exists within the on-chip budget, splitting the fusion is cheaper than inserting an HBM transpose, but the split must be explicit in the contract so its bytes are counted \citep{memoryfloor2026,taxbreak2026}. Pass ordering that runs layout legalization before fusion search prevents these negotiations from happening as emergency rewrites at codegen \citep{chen2020compilerbasedhardw,amd2024xdna}.

\paragraph{Fusion discovery order.}
Production pass pipelines should discover fusion in a fixed order: elementwise producers into their dominant consumers first, norm and scale operators into the GEMMs that follow, attention seams only after online-carrier legality is proven, and MoE routing last because it changes control flow \citep{dao2022flashattention,liu2024deepseekv3,chen2020compilerbasedhardw}. Each stage should emit a legality report that names the buffers preserved and the counters expected to move \citep{yirage2026,taxbreak2026,memoryfloor2026}. This ordering prevents a greedy pass from consuming a seam that later, more valuable fusion needs for carrier residency \citep{dao2022flashattention,patel2025llminference}. It also gives engineers a vocabulary for regressions: a fusion that disappears between releases can be traced to the stage that owns that seam \citep{yirage2026,chen2020compilerbasedhardw}.

\paragraph{Spill accounting discipline.}
Spills are not an error if they are counted; they are a bug when they are invisible \citep{memoryfloor2026,taxbreak2026,patel2025llminference}. Fusion passes should emit a spill event every time a planned on-chip buffer falls back to global memory, with the tensor name and byte width attached \citep{yirage2026,chen2020compilerbasedhardw}. Production dashboards should alert when spill bytes grow across compiler versions, because a silent unfusion and an over-eager new fusion produce the same spill signature \citep{memoryfloor2026,taxbreak2026}. Reviewers should ask fusion PRs to show the before-and-after spill ledger; a PR that reports only kernel count reductions has not yet demonstrated a bytes win \citep{patel2025llminference,semianalysis2024tma}.

\section{layout theory}
\label{sec:ch13_layout_theory}

\textbf{Layout theory} decides how tensors are stored and traversed: contiguous vs blocked layouts, KV cache page tables, head/interleaved formats, and alignment for TMA/vector loads \citep{kwon2023vllm,patel2025llminference,semianalysis2024tma,nvidia2022hopper}. A perfect fusion plan fails if layout forces strided HBM access---bytes/token rise while kernels look ``fused'' in the IR \citep{memoryfloor2026,yirage2026,taxbreak2026,chen2020compilerbasedhardw}.

PagedAttention (vLLM) is layout engineering for KV capacity: non-contiguous physical pages with logical continuity reduce fragmentation without changing model math \citep{kwon2023vllm,patel2025llminference,taxbreak2026}. Compilers must treat page table indirection as a first-class layout transform, not an runtime hack outside IR \citep{yirage2026,chen2020compilerbasedhardw,liu2024deepseekv3}. On NPU backends, layout often \emph{is} the schedule---static DMA descriptors replace pointer chasing \citep{amd2024xdna,amd2024aie,dlbook}.

\textbf{Phase-specific layouts.} Some teams store KV in a layout optimized for prefill writes (wide vector stores along sequence) and transposes for decode reads (head-local access) \citep{kwon2023vllm,dao2023flashdecoding,patel2025llminference}. Transpose costs bytes; the compiler must charge them in the cost model and prefer layouts that avoid per-token transposes when ctx grows \citep{memoryfloor2026,taxbreak2026,yirage2026,chen2020compilerbasedhardw,semianalysis2024tma,dlbook,amd2024xdna,amd2024aie,liu2024deepseekv3,dao2022flashattention,nvidia2022hopper,nvidia2024cudagraphs,dao2023flashdecoding}. When a layout pass cannot prove decode legality, split layouts per phase rather than forcing one format \citep{patel2025llminference,yirage2026,chen2020compilerbasedhardw,memoryfloor2026,taxbreak2026,dlbook}.

Layout passes interact with tiling: blocked weights pair with blocked GEMM tiles; NHWC vs NCHW choices change vectorization on CPU and GPU differently \citep{dlbook,chen2020compilerbasedhardw,semianalysis2024tma,patel2025llminference,memoryfloor2026,taxbreak2026,yirage2026,amd2024xdna,kwon2023vllm,liu2024deepseekv3,dao2022flashattention,dao2023flashdecoding,nvidia2022hopper,nvidia2024cudagraphs,amd2024aie}. Decode compilers should log layout transforms in pass metadata so regressions trace to reorder vs fusion vs tile changes \citep{yirage2026,chen2020compilerbasedhardw,patel2025llminference,taxbreak2026,memoryfloor2026,semianalysis2024tma,amd2024xdna,amd2024aie,dlbook,kwon2023vllm,liu2024deepseekv3,dao2022flashattention,dao2023flashdecoding,nvidia2022hopper,nvidia2024cudagraphs}.

\textbf{KV layout and paging.} Decode throughput scales with how cheaply each new token reads historical KV \citep{kwon2023vllm,dao2023flashdecoding,patel2025llminference}. Contiguous KV is simple but fragments GPU memory at long context; paged layouts add indirection cost that must be amortized by fewer reallocations and higher batch utilization \citep{kwon2023vllm,taxbreak2026}. Compiler layout passes should lower page-table updates into explicit IR ops so backends can fuse gather/scatter with attention tiles where legal \citep{yirage2026,chen2020compilerbasedhardw}.

\textbf{Weight layout and quantization.} INT4/FP8 weight layouts (block scales, interleaved scales) are layout transforms before they are ``quantization features'' \citep{liu2024deepseekv3,patel2025llminference,semianalysis2024tma}. A fusion plan that assumes FP16 contiguous weights will fail silently when weights are block-scaled---loads become gather-bound and bytes/token rise \citep{memoryfloor2026,yirage2026}. Layout legalization must run \emph{before} fusion search, not as a post-pass patch \citep{chen2020compilerbasedhardw,yirage2026,taxbreak2026}.

\textbf{Alignment and vector width.} TMA and vector units prefer aligned 128-byte lines; misaligned KV or activation tensors force uncoalesced access patterns that no tile size fixes \citep{semianalysis2024tma,nvidia2022hopper,dlbook}. CPU decode benefits from cache-line alignment across head dimensions; NPU backends often require fixed stride multiples baked into DMA descriptors \citep{amd2024xdna,amd2024aie,chen2020compilerbasedhardw}. Layout theory therefore includes alignment as a first-class constraint in IR, not a codegen afterthought \citep{yirage2026,patel2025llminference}.

\paragraph{Multi-hardware delta.} GPU layouts optimize for coalesced 128B lines; CPU layouts optimize for cache line reuse across batch-1 thin GEMMs; NPU layouts fix byte addresses at compile time \citep{nvidia2022hopper,amd2024xdna,dlbook,patel2025llminference,kwon2023vllm,yirage2026,chen2020compilerbasedhardw,taxbreak2026,memoryfloor2026,semianalysis2024tma,liu2024deepseekv3,dao2022flashattention,dao2023flashdecoding,nvidia2024cudagraphs,amd2024aie}.

\begin{example}
Switching KV layout to improve prefill vectorization can regress decode if batch-1 attention loads become gather-bound---layout theory requires phase-specific profiling \citep{kwon2023vllm,dao2023flashdecoding,patel2025llminference,taxbreak2026,memoryfloor2026,yirage2026,chen2020compilerbasedhardw,amd2024xdna,semianalysis2024tma,dlbook,liu2024deepseekv3,dao2022flashattention,nvidia2022hopper,nvidia2024cudagraphs,amd2024aie}.
\end{example}

Layout CI should snapshot tensor strides and alignment for KV, QKV weights, and MLP weights at each context bucket \citep{kwon2023vllm,yirage2026,chen2020compilerbasedhardw,patel2025llminference,semianalysis2024tma,liu2024deepseekv3,memoryfloor2026,taxbreak2026,dlbook,amd2024xdna,amd2024aie}. Stride drift without corresponding fusion changes is a common source of ``mysterious'' bytes/token regressions across compiler versions \citep{memoryfloor2026,taxbreak2026,yirage2026,chen2020compilerbasedhardw,patel2025llminference}.

\paragraph{Layout canonicalization and legality.}
Every phase should have a canonical layout derived from its access pattern, with explicit ops for any conversion rather than hidden runtime transposes \citep{semianalysis2024tma,yirage2026,chen2020compilerbasedhardw}. Legality in layout terms means: the chosen strides are aligned for the target's vector or DMA width, the tensor fits its assigned residency class, and no consumer requires an implicit reshape \citep{nvidia2022hopper,amd2024aie,patel2025llminference}. A layout pass that cannot prove these three properties should refuse to emit rather than insert a conversion that doubles bytes/token \citep{memoryfloor2026,taxbreak2026}. CI should snapshot the canonical layout of KV, QKV weights, and MLP weights per context bucket so drift appears as a first-class regression \citep{chen2020compilerbasedhardw,dlbook}.

\paragraph{Paged KV layout decisions.}
Paging changes the layout question from ``where is the next KV row'' to ``which page does the block walk touch,'' and the answer determines whether attention loads coalesce \citep{kwon2023vllm,dao2023flashdecoding,patel2025llminference}. Page size should be chosen against the head dimension and the on-chip staging width: pages that are too small scatter the block walk, while pages that are too large waste capacity on partial occupancy \citep{semianalysis2024tma,memoryfloor2026}. The compiler should expose the page-table walk as an explicit IR pattern so fusion can keep gather and attention in one region where the target allows it \citep{yirage2026,chen2020compilerbasedhardw}. On XDNA, page indirection may need to become a static descriptor list because runtime pointer chasing is not a legal DMA pattern \citep{amd2024xdna,amd2024aie}.

\paragraph{Weight layout conversion timing.}
Quantized and block-scaled weights should be converted once at load or export time, never per decode step, and the conversion cost must appear in the same ledger as activation traffic \citep{liu2024deepseekv3,patel2025llminference,memoryfloor2026}. A compiler that treats weight conversion as a free pre-pass hides a real cost when models are swapped at runtime or when context buckets share one weight copy \citep{chen2020compilerbasedhardw,yirage2026}. Layout legalization should verify that every quantized tensor has a single canonical runtime layout and that no fusion plan inserts a second conversion \citep{patel2025llminference,chen2020compilerbasedhardw}. Runtime weight swaps are where layout bugs surface as mysteriously slow first tokens; the fix belongs in the pass pipeline, not in cache tuning \citep{taxbreak2026,memoryfloor2026}.

\section{parallel scheduling}
\label{sec:ch13_parallel_scheduling}

\textbf{Parallel scheduling} maps work to SMs, cores, or AIE arrays---including Flash-Decoding's parallelization over KV sequence length at batch-1 \citep{dao2023flashdecoding,patel2025llminference,taxbreak2026}. Prefill parallelizes over tokens; decode parallelizes over heads, KV blocks, or expert routes while issuing one new token per step \citep{kwon2023vllm,dlbook,dao2022flashattention,memoryfloor2026,yirage2026,chen2020compilerbasedhardw,semianalysis2024tma,amd2024xdna,liu2024deepseekv3,dao2023flashdecoding,nvidia2022hopper,nvidia2024cudagraphs,amd2024aie,taxbreak2026,patel2025llminference}.

The failure mode is ``parallelism for idle hardware'': launching extra blocks that contend for HBM bandwidth without increasing useful bytes delivered per ms/token \citep{memoryfloor2026,semianalysis2024tma,patel2025llminference,taxbreak2026}. Schedulers should be bandwidth-aware---when $R_{\mathrm{floor}}$ is high, additional parallel splits hurt \citep{memoryfloor2026,taxbreak2026,yirage2026,chen2020compilerbasedhardw,amd2024xdna,amd2024aie,dlbook,kwon2023vllm,liu2024deepseekv3,dao2022flashattention,dao2023flashdecoding,nvidia2022hopper,nvidia2024cudagraphs,semianalysis2024tma,patel2025llminference}. When launches/token dominate (TaxBreak regime), parallel capture inside CUDA Graphs or fused kernels may help even at moderate bandwidth \citep{taxbreak2026,nvidia2024cudagraphs,memoryfloor2026,yirage2026,chen2020compilerbasedhardw,patel2025llminference,amd2024xdna,semianalysis2024tma,dlbook,kwon2023vllm,liu2024deepseekv3,dao2022flashattention,dao2023flashdecoding,nvidia2022hopper,amd2024aie}.

\textbf{Prefill/decode scheduler split.} Prefill schedulers maximize token-parallel work; decode schedulers maximize memory efficiency per new token \citep{patel2025llminference,dao2022flashattention,dao2023flashdecoding,kwon2023vllm,taxbreak2026,memoryfloor2026,yirage2026,chen2020compilerbasedhardw,dlbook,amd2024xdna,amd2024aie,semianalysis2024tma,liu2024deepseekv3,nvidia2022hopper,nvidia2024cudagraphs,dao2023flashdecoding}. Sharing one scheduler heuristic across phases is a common source of decode regressions after prefill optimizations \citep{patel2025llminference,memoryfloor2026,taxbreak2026,yirage2026,chen2020compilerbasedhardw,dlbook}.

Compiler IR should expose parallel axes (batch, head, KV block, expert) so backends legalize splits per target \citep{yirage2026,chen2020compilerbasedhardw,liu2024deepseekv3,amd2024aie,amd2024xdna,nvidia2022hopper,semianalysis2024tma,patel2025llminference,taxbreak2026,memoryfloor2026,kwon2023vllm,dlbook,dao2022flashattention,dao2023flashdecoding,nvidia2024cudagraphs}. Static NPU schedules often fix parallelism at compile time; GPU schedules retain dynamic grid sizing within capture buckets \citep{amd2024xdna,nvidia2024cudagraphs,chen2020compilerbasedhardw,yirage2026,patel2025llminference,taxbreak2026,memoryfloor2026,semianalysis2024tma,dlbook,kwon2023vllm,liu2024deepseekv3,dao2022flashattention,dao2023flashdecoding,nvidia2022hopper,amd2024aie}.

\textbf{Occupancy vs bandwidth.} GPU textbooks emphasize occupancy; decode compilers must ask whether extra warps increase \emph{useful} bytes delivered per cycle or merely contend for the same HBM port \citep{memoryfloor2026,semianalysis2024tma,patel2025llminference,taxbreak2026}. When $R_{\mathrm{floor}}$ is near unity, the scheduler should prefer fewer, wider memory transactions (better fusion, better layout) over more parallel blocks \citep{memoryfloor2026,dao2022flashattention,yirage2026}.

\textbf{Wave quantization and tail effects.} Parallel splits that do not divide sequence length evenly leave tail waves with low utilization---common when KV blocks do not align with context length \citep{dao2023flashdecoding,kwon2023vllm,patel2025llminference}. Compiler scheduling should pad or bucket context lengths explicitly so tail penalties appear in CI, not only at customer ctx${=}8192$ \citep{memoryfloor2026,taxbreak2026,yirage2026,chen2020compilerbasedhardw}.

\textbf{MoE and expert parallelism.} MoE layers introduce a routing axis: experts parallelize differently from attention heads \citep{liu2024deepseekv3,taxbreak2026,kwon2023vllm}. Schedulers must avoid launching all experts when only top-$k$ activate---parallel theory includes \emph{conditional} parallelism gated on router output, which complicates graph capture and static NPU schedules alike \citep{taxbreak2026,nvidia2024cudagraphs,amd2024aie,yirage2026,chen2020compilerbasedhardw,patel2025llminference}.

\paragraph{Multi-hardware delta.} GPU parallel scheduling uses grids/warps; CPU uses thread pools with cache affinity; NPU uses spatial arrays with fixed slot counts \citep{nvidia2022hopper,amd2024xdna,dlbook,chen2020compilerbasedhardw,yirage2026,patel2025llminference,taxbreak2026,memoryfloor2026,kwon2023vllm,liu2024deepseekv3,dao2022flashattention,semianalysis2024tma,dao2023flashdecoding,nvidia2024cudagraphs,amd2024aie}.

\begin{example}
Flash-Decoding reports large speedups at long context by parallelizing KV---the scheduler must still respect on-chip softmax carrier residency from Chapter~10 or parallel blocks reintroduce HBM traffic \citep{dao2023flashdecoding,dao2022flashattention,yirage2026,memoryfloor2026,taxbreak2026,patel2025llminference,chen2020compilerbasedhardw,amd2024xdna,semianalysis2024tma,dlbook,kwon2023vllm,liu2024deepseekv3,nvidia2022hopper,nvidia2024cudagraphs,amd2024aie}.
\end{example}

Parallel scheduling reviews should attach Nsight ``memory throughput'' and ``SM busy'' together---high SM busy with flat bytes/token indicates useless parallelism \citep{memoryfloor2026,semianalysis2024tma,patel2025llminference,taxbreak2026,yirage2026,chen2020compilerbasedhardw,dao2022flashattention,dao2023flashdecoding,dlbook,amd2024xdna,amd2024aie,kwon2023vllm,liu2024deepseekv3,nvidia2022hopper,nvidia2024cudagraphs}. Decode schedulers that only maximize occupancy violate the bandwidth-aware principle this section states \citep{memoryfloor2026,taxbreak2026,patel2025llminference}.

\paragraph{Ranking parallel axes for decode.}
Decode exposes four parallel axes---batch rows, query heads, KV blocks, and expert routes---and the scheduler should rank them by what each axis costs in merged state \citep{dao2023flashdecoding,kwon2023vllm,liu2024deepseekv3}. Parallelizing over query heads is cheap until shared scores or shared KV demand a merge; parallelizing over KV blocks always ends in an online-softmax merge; parallelizing over batch rows is the only axis that adds independent work without merge traffic \citep{dao2022flashattention,patel2025llminference}. At batch-1, batch rows vanish, which pushes decode toward head and KV parallelism and makes the merge tree a first-class scheduling cost \citep{taxbreak2026,memoryfloor2026}. The axis ranking should be recomputed per SKU because merge costs are residency costs, and residency tables differ by target \citep{amd2024xdna,yirage2026,chen2020compilerbasedhardw}.

\paragraph{Bandwidth-saturation test for any split.}
Before merging a parallelization change, run the saturation test: double the split count while holding the algorithm fixed and observe whether bytes/token stays flat \citep{memoryfloor2026,semianalysis2024tma,taxbreak2026}. If doubling splits raises bytes/token, the extra blocks are paying for their own reloads; if ms/token drops while bytes stay flat, the added parallelism is hiding latency productively \citep{patel2025llminference,memoryfloor2026}. The same test applies to capture groups inside graphs: a graph that launches more blocks than the memory system can feed merely serializes behind bandwidth \citep{nvidia2024cudagraphs,taxbreak2026}. Publishing the saturation curve with each kernel gives reviewers an objective reason to accept or reject a split count \citep{yirage2026,chen2020compilerbasedhardw}.

\paragraph{Static versus dynamic parallelism.}
GPU schedulers can choose grid sizes at launch within capture buckets, while XDNA/AIE fixes the spatial mapping at compile time and CPU leaves parallelism to the runtime thread pool \citep{nvidia2024cudagraphs,amd2024aie,dlbook}. Compiler IR should annotate which axes are static and which are dynamic so the same parallel plan can be legalized per backend \citep{yirage2026,chen2020compilerbasedhardw}. A dynamic axis on GPU that becomes static on NPU changes tail behavior: the static schedule must include worst-case splits because it cannot adapt to short context \citep{amd2024xdna,amd2024aie}. Bucketing context lengths is therefore not just a graph-capture concern; it is the mechanism that keeps static schedules from wasting array capacity on tail waves \citep{kwon2023vllm,taxbreak2026}.

\paragraph{Merge cost as a schedule term.}
Every split that produces partial state must pay a merge, and the schedule should treat merge bytes and merge latency as explicit costs rather than epilogue noise \citep{dao2023flashdecoding,taxbreak2026,patel2025llminference}. When split count doubles, merge state roughly doubles with it, so the bandwidth-saturation test must measure the merge stage separately from the compute stage \citep{memoryfloor2026,dao2022flashattention}. Compilers should annotate split boundaries with their merge tree so autotuners can compare one large split against two smaller ones on equal terms \citep{yirage2026,chen2020compilerbasedhardw}. A merge that stays on-chip is nearly free; a merge that touches global memory can negate the split's entire benefit \citep{semianalysis2024tma,memoryfloor2026}.

\section{GPU CPU NPU split}
\label{sec:ch13_gpu_cpu_npu_split}

\textbf{GPU CPU NPU split} is how one decoder IR lowers to CUDA/HIP, CPU oneDNN/OpenMP paths, and XDNA/AIE static dataflow without rewriting model Python \citep{yirage2026,chen2020compilerbasedhardw,amd2024xdna,amd2024aie,patel2025llminference,taxbreak2026,memoryfloor2026,dlbook,kwon2023vllm,liu2024deepseekv3,dao2022flashattention,semianalysis2024tma,dao2023flashdecoding,nvidia2022hopper,nvidia2024cudagraphs}. Heterogeneous fleets fail when the GPU schedule is copied to NPU---illegal DMA edges and tile overflows appear as silent CPU fallbacks or wrong results \citep{amd2024xdna,amd2024aie,yirage2026,chen2020compilerbasedhardw,patel2025llminference,taxbreak2026,memoryfloor2026,dlbook,kwon2023vllm,liu2024deepseekv3,dao2022flashattention,semianalysis2024tma,dao2023flashdecoding,nvidia2022hopper,nvidia2024cudagraphs}.

Backend split should share high-level passes (fusion legality, layout, tiling metadata) and fork only at emit \citep{yirage2026,chen2020compilerbasedhardw,liu2024deepseekv3,amd2024xdna,amd2024aie,nvidia2022hopper,semianalysis2024tma,patel2025llminference,taxbreak2026,memoryfloor2026,kwon2023vllm,dlbook,dao2022flashattention,dao2023flashdecoding,nvidia2024cudagraphs}. Triplet regression---same model, same context buckets, GPU vs CPU vs NPU---is the acceptance test for split correctness \citep{patel2025llminference,yirage2026,chen2020compilerbasedhardw,amd2024xdna,amd2024aie,taxbreak2026,memoryfloor2026,semianalysis2024tma,dlbook,kwon2023vllm,liu2024deepseekv3,dao2022flashattention,dao2023flashdecoding,nvidia2022hopper,nvidia2024cudagraphs}. Metrics remain bytes/token and launches/token per target; FLOPs parity is irrelevant \citep{taxbreak2026,memoryfloor2026,patel2025llminference,yirage2026,chen2020compilerbasedhardw,amd2024xdna,amd2024aie,dlbook,kwon2023vllm,liu2024deepseekv3,dao2022flashattention,semianalysis2024tma,dao2023flashdecoding,nvidia2022hopper,nvidia2024cudagraphs}.

\textbf{Version skew.} GPU and NPU toolchains release on different cadences; shared IR must pin legality rules so a middle-end upgrade cannot silently break NPU emit while GPU CI stays green \citep{yirage2026,chen2020compilerbasedhardw,amd2024aie,amd2024xdna,patel2025llminference,taxbreak2026,memoryfloor2026,dlbook,kwon2023vllm,liu2024deepseekv3,dao2022flashattention,semianalysis2024tma,dao2023flashdecoding,nvidia2022hopper,nvidia2024cudagraphs}.

CPU backends often serve hybrid pipelines: dynamic shapes, sampling, and small ops stay on host while GPU/NPU run fused decode cores \citep{kwon2023vllm,patel2025llminference,yirage2026,chen2020compilerbasedhardw,taxbreak2026,memoryfloor2026,amd2024xdna,amd2024aie,dlbook,liu2024deepseekv3,dao2022flashattention,semianalysis2024tma,dao2023flashdecoding,nvidia2022hopper,nvidia2024cudagraphs}. NPU backends require upfront static shape buckets; GPU backends tolerate bucketed graph capture \citep{nvidia2024cudagraphs,amd2024aie,yirage2026,chen2020compilerbasedhardw,patel2025llminference,taxbreak2026,memoryfloor2026,amd2024xdna,semianalysis2024tma,dlbook,kwon2023vllm,liu2024deepseekv3,dao2022flashattention,dao2023flashdecoding,nvidia2022hopper}. Document which ops each backend owns so mode selection (Chapter~12) maps cleanly to emit paths \citep{yirage2026,chen2020compilerbasedhardw,patel2025llminference,taxbreak2026,memoryfloor2026,nvidia2024cudagraphs,amd2024xdna,amd2024aie,semianalysis2024tma,dlbook,kwon2023vllm,liu2024deepseekv3,dao2022flashattention,dao2023flashdecoding,nvidia2022hopper}.

\textbf{Shared middle-end, forked back-end.} The durable pattern is a shared middle-end (fusion/layout/tiling metadata) and forked backends that interpret residency classes per target \citep{yirage2026,chen2020compilerbasedhardw,amd2024xdna,amd2024aie,dlbook}. Glow and TVM historically split early; LLM decode stacks should delay fork until after decode-specific legality checks \citep{chen2020compilerbasedhardw,patel2025llminference,yirage2026,taxbreak2026}.

\textbf{Fallback paths.} When NPU legality fails, frameworks silently fall back to CPU loops---correct but catastrophically slow \citep{amd2024xdna,amd2024aie,yirage2026,patel2025llminference}. Production compilers must treat fallback as an explicit compile error unless a documented degraded mode is enabled \citep{chen2020compilerbasedhardw,yirage2026,memoryfloor2026,taxbreak2026}. Triplet CI should fail when GPU bytes/token improve but NPU emits empty kernels \citep{patel2025llminference,amd2024aie,amd2024xdna}.

\textbf{Heterogeneous single-node pipelines.} Some deployments run attention on GPU and MLP on NPU, or prefill on GPU and decode on CPU \citep{patel2025llminference,liu2024deepseekv3,kwon2023vllm}. Split theory then includes \emph{partition} passes that cut the graph at device boundaries with explicit copy ops whose bytes count toward goodput \citep{yirage2026,chen2020compilerbasedhardw,taxbreak2026,memoryfloor2026,dlbook}. Hidden host-device copies dominate when partition points are chosen for convenience rather than residency \citep{patel2025llminference,taxbreak2026}.

\paragraph{Multi-hardware delta.} GPU emphasizes TMA+MMA and graphs; CPU emphasizes cache blocking and NUMA; NPU emphasizes static DMA+tiles---same IR, three residency models \citep{nvidia2022hopper,amd2024xdna,dlbook,yirage2026,chen2020compilerbasedhardw,patel2025llminference,taxbreak2026,memoryfloor2026,kwon2023vllm,liu2024deepseekv3,dao2022flashattention,semianalysis2024tma,dao2023flashdecoding,nvidia2024cudagraphs,amd2024aie}.

\begin{example}
A fusion plan legal on Hopper shared memory may be rejected on XDNA---split backends must surface legality errors at compile time, not as empty NPU binaries in production \citep{yirage2026,amd2024xdna,amd2024aie,chen2020compilerbasedhardw,patel2025llminference,taxbreak2026,memoryfloor2026,dlbook,kwon2023vllm,liu2024deepseekv3,dao2022flashattention,semianalysis2024tma,dao2023flashdecoding,nvidia2022hopper,nvidia2024cudagraphs}.
\end{example}

Release gates for heterogeneous builds should require artifact parity: every target emits a non-empty binary, passes numerics smoke tests, and reports bytes/token within agreed triplet tolerance \citep{yirage2026,chen2020compilerbasedhardw,patel2025llminference,amd2024xdna,amd2024aie,taxbreak2026,memoryfloor2026,dlbook,kwon2023vllm,liu2024deepseekv3,dao2022flashattention,semianalysis2024tma,dao2023flashdecoding,nvidia2022hopper,nvidia2024cudagraphs}. GPU-only green CI is insufficient for books and products that claim cross-hardware compilation \citep{patel2025llminference,yirage2026,chen2020compilerbasedhardw,amd2024aie,amd2024xdna,dlbook}.

\begin{table}[t]
\centering
\caption{Core theory pillars mapped to decode counters (profile per SKU).}
\label{tab:ch13_pillars}
\begin{tabular}{p{0.18\linewidth}p{0.38\linewidth}p{0.32\linewidth}}
\toprule
\textbf{Pillar} & \textbf{Compiler question} & \textbf{Primary counter} \\
\midrule
Tiling & Which blocks stay on-chip? & bytes/token \\
Fusion & Which boundaries can vanish? & bytes/token, launches/token \\
Layout & How are KV/weights traversed? & bytes/token \\
Parallel & Which axes saturate bandwidth? & ms/token, SM util. \\
GPU/CPU/NPU split & Where does each op emit? & triplet parity \\
\bottomrule
\end{tabular}
\end{table}

\paragraph{Operator ownership matrix.}
Backend splits are easier to maintain when the team publishes an ownership matrix: which operators emit natively on GPU, CPU, and NPU, which emit through a shared library path, and which are explicitly unsupported on each target \citep{yirage2026,chen2020compilerbasedhardw,amd2024xdna}. Empty cells in the matrix are legal only when a documented fallback exists; an undocumented empty cell is how fleet incidents begin with a silent CPU loop \citep{amd2024aie,patel2025llminference,taxbreak2026}. The matrix should be generated from the IR test suite so it cannot drift from what the backends actually accept \citep{memoryfloor2026,chen2020compilerbasedhardw}. Ownership disputes---for example attention fused by the GPU backend but rejected by the NPU backend---should surface as a legality report at the shared pass level, not as a fork in model code \citep{yirage2026,amd2024aie}.

\paragraph{Partition copy cost model.}
Heterogeneous pipelines that cut the graph at device boundaries need a cost model for the copies themselves: host-device and device-device transfers consume bandwidth and add latency, and a partition chosen for API convenience can cost more than the kernel it isolates \citep{patel2025llminference,taxbreak2026,memoryfloor2026}. The partition pass should evaluate each candidate cut with the copy bytes charged against the same bytes/token ledger used for fusion \citep{chen2020compilerbasedhardw,yirage2026}. When a copy is unavoidable, prefer cuts at tensors that are small relative to the compute they enable, and fuse the transfer with the following kernel where the hardware supports it \citep{semianalysis2024tma,amd2024xdna}. Release gates should show the copy bytes explicitly; hiding them behind profiler rows is how heterogeneous deployments look fused and stay slow \citep{patel2025llminference,memoryfloor2026}.

\paragraph{Pass ordering and golden IR diffs.}
The five pillars interact through ordering, so the pass pipeline itself should be versioned and tested as a unit \citep{chen2020compilerbasedhardw,yirage2026}. Each release should capture a golden IR diff showing the sequence tiling classification, layout canonicalization, fusion contracts, parallel annotations, and backend splits \citep{memoryfloor2026,taxbreak2026}. A golden diff turns subtle ordering changes into reviewable artifacts rather than mysterious fleet regressions \citep{patel2025llminference,chen2020compilerbasedhardw}. When a new pass is proposed, the review question is not whether it is clever but where it sits in the ordering and which counter it moves on batch-1 decode cells \citep{yirage2026,taxbreak2026,memoryfloor2026}.

\paragraph{Triplet tolerance policy.}
Cross-target parity needs an explicit tolerance policy: which counters must match within a percentage, which may legitimately differ because residency classes differ, and which mismatches are always failures \citep{patel2025llminference,yirage2026,chen2020compilerbasedhardw}. Bytes/token and launches/token should be compared per target against that target's own baseline rather than against a single global number, because an NPU that eliminates launches entirely has a different launch definition than a GPU \citep{amd2024xdna,amd2024aie,memoryfloor2026}. Numerics tolerance should be defined at the attention-output level, not per op, so backend reordering that stays within output tolerance is allowed \citep{dao2022flashattention,patel2025llminference}. Without a written policy, triplet CI becomes a source of noise or, worse, is skipped when numbers disagree \citep{taxbreak2026,memoryfloor2026}.

\section{Key Takeaways}
\label{sec:ch13_key_takeaways}

\begin{itemize}
\item \textbf{Tiling is IO theory, not FLOP theory.}
FlashAttention-style blocks exist to choose which sub-blocks of weights, activations, and KV fit in registers, shared memory, L2, or NPU tile SRAM for each decode step, so the goal is cutting HBM traffic \citep{dao2022flashattention,semianalysis2024tma}. Tile search must therefore be driven by decode bytes/token, not by GEMM FLOPs, or it optimizes the wrong objective \citep{memoryfloor2026,taxbreak2026}.

\item \textbf{Fusion has legality, not just opportunity.}
Each unfused operator boundary exports activations to HBM, but greedy fusion that ignores carrier and residency checks silently regresses decode by overflowing the on-chip budget \citep{dao2022flashattention,yirage2026}. Fusion theory governs \emph{when} a boundary may legally disappear, which a compiler must verify before merging operators \citep{chen2020compilerbasedhardw,patel2025llminference}.

\item \textbf{Layout is a compiler transform, not a runtime afterthought.}
How tensors are stored and traversed---contiguous versus blocked, paged KV tables, head-interleaved formats, TMA/vector alignment---decides whether a fusion is even reachable \citep{kwon2023vllm,patel2025llminference}. A wrong layout can double bytes/token without changing FLOPs, so layout assignment belongs in the pass pipeline ahead of fusion \citep{taxbreak2026,yirage2026}.

\item \textbf{Parallelism must be bandwidth-aware.}
Mapping work to SMs, cores, or AIE arrays---including Flash-Decoding's parallelization over KV length at batch-1---only helps when on-chip residency holds across the split \citep{dao2023flashdecoding,semianalysis2024tma}. A parallel axis that forces activations off-chip trades latency hiding for bandwidth it cannot afford \citep{memoryfloor2026,patel2025llminference}.

\item \textbf{One IR, split backends, mandatory triplet regression.}
A single decoder IR lowers to CUDA/HIP, CPU oneDNN/OpenMP, and XDNA/AIE static dataflow with genuinely different emit, so correctness on one target says nothing about the others \citep{yirage2026,amd2024xdna,amd2024aie}. The five pillars are co-dependent---a layout change can illegalize a fusion, a fusion can force new tiles, a parallel split can break capture---so develop them as co-design profiled on decode counters, gated in CI by golden IR diffs and triplet bytes/token \citep{chen2020compilerbasedhardw,memoryfloor2026,taxbreak2026}.
\end{itemize}

The pillars are not independent knobs---a layout change can illegalize a fusion plan; a fusion plan can force new tile sizes; a parallel split can break graph capture \citep{yirage2026,chen2020compilerbasedhardw,memoryfloor2026,nvidia2024cudagraphs,patel2025llminference,taxbreak2026,dao2022flashattention,semianalysis2024tma,amd2024xdna,amd2024aie,dlbook,kwon2023vllm,liu2024deepseekv3,dao2023flashdecoding,nvidia2022hopper}. Treat compiler development as co-design across all five, profiled on decode counters, not as a sequence of isolated dialect tickets \citep{patel2025llminference,taxbreak2026,memoryfloor2026,yirage2026,chen2020compilerbasedhardw,dlbook}.

\section{Conclusion}
\label{sec:ch13_conclusion}

Core compiler theory---tiling, fusion, layout, parallel scheduling, and heterogeneous split---is the machinery behind Chapter~12 execution modes and the MegaKernel chapters that follow \citep{yirage2026,chen2020compilerbasedhardw,taxbreak2026,memoryfloor2026,patel2025llminference,amd2024xdna,amd2024aie,dao2022flashattention,dao2023flashdecoding,kwon2023vllm,liu2024deepseekv3,nvidia2022hopper,semianalysis2024tma,nvidia2024cudagraphs,dlbook}. Pass ordering slides are insufficient; profiled decode counters justify every transform \citep{patel2025llminference,taxbreak2026,memoryfloor2026,semianalysis2024tma,yirage2026,chen2020compilerbasedhardw,amd2024xdna,amd2024aie,dao2022flashattention,dao2023flashdecoding,kwon2023vllm,liu2024deepseekv3,nvidia2022hopper,nvidia2024cudagraphs,dlbook}.

Reviewers should ask five questions before approving any compiler PR: (1) Which pillar moves? (2) Which counter must improve on BS${=}1$ decode? (3) Which Chapter~10--11 constraint applies? (4) Does layout legalize before fusion? (5) Does triplet regression pass on GPU, CPU, and NPU? \citep{yirage2026,chen2020compilerbasedhardw,patel2025llminference,taxbreak2026,memoryfloor2026,amd2024xdna,amd2024aie,dlbook}. If any answer is ``unknown,'' defer the patch until telemetry exists \citep{patel2025llminference,taxbreak2026,memoryfloor2026}. These questions take minutes in review and prevent weeks of fleet debugging on batch-1 decode cells at production context buckets in CI gates \citep{yirage2026,chen2020compilerbasedhardw,patel2025llminference,memoryfloor2026,taxbreak2026,dlbook}.

Chapter~14 instantiates these principles in MLIR dialects and pass pipelines---where tiling, fusion, and layout become rewrite patterns you can test in CI \citep{chen2020compilerbasedhardw,yirage2026,liu2024deepseekv3,patel2025llminference,taxbreak2026,memoryfloor2026,amd2024xdna,amd2024aie,semianalysis2024tma,dao2022flashattention,dao2023flashdecoding,kwon2023vllm,nvidia2022hopper,nvidia2024cudagraphs,dlbook}. Carry Table~\ref{tab:ch13_pillars} forward: if a dialect change cannot name which pillar moves which counter, defer the patch \citep{yirage2026,chen2020compilerbasedhardw,patel2025llminference,taxbreak2026,memoryfloor2026,amd2024xdna,amd2024aie,dlbook,kwon2023vllm,liu2024deepseekv3,dao2022flashattention,semianalysis2024tma,dao2023flashdecoding,nvidia2022hopper,nvidia2024cudagraphs}.

Until then, run one decode profile per proposed pass---BS${=}1$, ctx buckets $\{512,2048,8192\}$, bytes/token and launches/token on GPU and at least one non-GPU target \citep{patel2025llminference,memoryfloor2026,taxbreak2026,yirage2026,chen2020compilerbasedhardw,amd2024xdna,amd2024aie,dlbook,kwon2023vllm,liu2024deepseekv3,dao2022flashattention,semianalysis2024tma,dao2023flashdecoding,nvidia2022hopper,nvidia2024cudagraphs}. That habit converts abstract compiler theory into the measurable goodput improvements Part~VI promises \citep{taxbreak2026,memoryfloor2026,patel2025llminference,yirage2026,chen2020compilerbasedhardw,dlbook}. Part~VI continues with dialect-level mechanics; carry these pillars as the acceptance language for every rewrite pattern \citep{chen2020compilerbasedhardw,yirage2026,liu2024deepseekv3,patel2025llminference,taxbreak2026,memoryfloor2026,amd2024xdna,amd2024aie,semianalysis2024tma,dlbook,kwon2023vllm,dao2022flashattention,dao2023flashdecoding,nvidia2022hopper,nvidia2024cudagraphs}.

\typeout{END_CHAPTER "ch13" \theabspage}
